\chapter{Determinación estructural: RMN de \texorpdfstring{$^{13}$C}{13C}, EM e IR combinadas}\label{ch:b2:structure-determination}

Un laboratorio de perfumería recibe un vial de un líquido incoloro que huele a jazmín. Nadie ha
escrito una etiqueta. Un \omterm{def:b2:mass-spec-atomic:spectrum}{espectro de masas} da su fórmula molecular, un espectro infrarrojo sus grupos
funcionales, un espectro de carbono 13 su esqueleto y un espectro de protón la manera en que se unen las piezas; en una
hora el líquido tiene nombre. Este capítulo añade la \omterm{def:b2:structure-determination:c13}{RMN de carbono 13} a la RMN de protón y a la espectroscopia infrarroja
del volumen del primer año y a la \omterm{def:b2:mass-spec-atomic:spectrum}{espectrometría de masas} del \cref{ch:b2:mass-spec-atomic}, y convierte las cuatro
en un solo método.

\begin{recall}
El volumen del primer año: desplazamiento químico y apantallamiento, protones equivalentes, integración, acoplamiento espín--espín
y regla $n + 1$, números de onda infrarrojos de los grupos funcionales y región de la huella dactilar, grado
de insaturación, resolución de una estructura a partir de tres espectros. \Cref{ch:b2:mass-spec-atomic}: el \omterm{def:b2:mass-spec-atomic:molecular-ion}{ion
molecular}, los \omterm{def:b2:mass-spec-atomic:isotope-pattern}{perfiles isotópicos}, la fragmentación.
\end{recall}

\section{RMN de carbono 13}

El carbono 12, el isótopo principal, no tiene espín nuclear y no da señal de RMN. El carbono 13 tiene un espín de
$\tfrac12$, como el protón, pero solo constituye el 1.1~\% del carbono natural y su señal es más débil,
núcleo por núcleo, que la de un protón; un espectro de carbono necesita por ello más muestra o más acumulaciones.
% ledger: iso:C13

\begin{definition}[RMN de carbono 13]\label{def:b2:structure-determination:c13}
La \emph{RMN de carbono 13}\index{RMN de carbono 13} registra las resonancias de los núcleos \ce{^{13}C} de una muestra.
Suele realizarse con \emph{desacoplamiento de banda ancha}\index{desacoplamiento de banda ancha}: los protones se
irradian en todo su intervalo de frecuencias durante la medida, lo que elimina todo
acoplamiento carbono--protón, de modo que cada tipo de carbono da una sola línea. Los \emph{carbonos
equivalentes}\index{carbonos equivalentes}, intercambiados por una simetría de la molécula (o por una rotación rápida),
dan la misma línea.
\end{definition}

\begin{proposition}[Sin acoplamiento carbono--carbono]\label{prop:b2:structure-determination:no-cc-coupling}
Los acoplamientos entre átomos de carbono vecinos no aparecen en un espectro de \ce{^{13}C} en abundancia
natural.
\end{proposition}

\begin{proof}
Un acoplamiento entre dos carbonos solo se observa en las moléculas en las que ambos son \ce{^{13}C}. Para dos posiciones
vecinas dadas, la probabilidad es $a_{13}^2 = 0.011^2 \approx 1.2 \times 10^{-4}$: alrededor de una
molécula de cada ocho mil. La señal de cada carbono procede casi enteramente de moléculas en las que sus
vecinos son \ce{^{12}C}, invisibles para el espectrómetro; los satélites acoplados son cien veces
más débiles que el ruido de un espectro ordinario.
\end{proof}

\begin{proposition}[Número de líneas]\label{prop:b2:structure-determination:signals}
Un espectro de \ce{^{13}C} desacoplado muestra una línea por cada conjunto de \omterm{def:b2:structure-determination:c13}{carbonos equivalentes}, salvo que dos conjuntos
tengan por casualidad el mismo desplazamiento.
\end{proposition}

\begin{proof}[Argumento]
Los \omterm{def:b2:structure-determination:c13}{carbonos equivalentes} están en entornos idénticos, de modo que tienen el mismo desplazamiento; los carbonos no equivalentes
suelen diferir. Eliminados los acoplamientos con los protones y ausentes los acoplamientos entre carbonos
(\cref{prop:b2:structure-determination:no-cc-coupling}), nada desdobla una línea. El solapamiento accidental
se produce cuando dos carbonos distintos tienen desplazamientos más próximos que la resolución espectral, lo que es frecuente para los
carbonos de un anillo bencénico hacia \qty{128}{ppm}.
\end{proof}

Los desplazamientos se extienden sobre unos \qty{220}{ppm}, veinte veces el intervalo de los protones, de modo que los solapamientos son más raros que
en los espectros de protón, y el desplazamiento por sí solo indica el tipo de carbono. El gráfico siguiente muestra los desplazamientos
medidos para los compuestos de este capítulo.

\begin{omfigure}
\begin{tikzpicture}
\begin{axis}[width=0.86\linewidth, height=4.6cm, xmin=0, xmax=220, ymin=-0.6, ymax=4.6, x dir=reverse,
  xlabel={$\delta$ (ppm)}, label style={font=\small}, tick label style={font=\scriptsize},
  ytick={0,1,2,3,4}, yticklabels={C alquílico, C--O, C aromático, C=O de éster, C=O de cetona},
  xtick={0,20,40,60,80,100,120,140,160,180,200,220}, grid=major, grid style={black!10}]
\addplot[only marks, mark=|, mark size=5pt, thick, omDef] table[x=shift, y=cls] {figdata/out/bachelor-2-nmr-spectra-d.dat};
\end{axis}
\end{tikzpicture}

{\small Desplazamientos de \ce{^{13}C} medidos para la butan-2-ona, la pentan-2-ona, el benzoato de etilo y el etanoato de bencilo,
ordenados por tipo de carbono (CDCl$_3$; datos de PubChem). Los carbonilos de cetona se sitúan hacia 209, los de éster hacia
167--171, los carbonos \omterm{def:b2:huckel:aromatic}{aromáticos} entre 128 y 137, los carbonos unidos a un oxígeno de éster hacia 61--66, y los
carbonos alquílicos por debajo de 46.}
% figdata: bachelor-2/nmr-spectra.py part d
% ledger: c13:butanone, c13:pentan-2-one, c13:ethylbenzoate, c13:benzylacetate
\end{omfigure}

\begin{proposition}[Las alturas no son recuentos]\label{prop:b2:structure-determination:intensities}
En un espectro de \ce{^{13}C} desacoplado de rutina, las alturas de las líneas no son proporcionales a los
números de carbonos que representan; los carbonos sin hidrógenos dan líneas débiles.
\end{proposition}

\begin{proof}[Argumento]
Los espectros de rutina repiten la medida más deprisa de lo que los carbonos más lentos tardan en relajarse al equilibrio, y
los carbonos sin protones unidos se relajan lentamente, de modo que su señal está parcialmente saturada. El desacoplamiento además
refuerza las señales de los carbonos que llevan protones, en una cantidad que depende de su entorno.
En la butan-2-ona, la línea del carbonilo es la más débil de las cuatro, aunque representa un carbono como
las demás. La integración, tan útil para los protones, no se usa por tanto en los espectros de carbono de rutina.
\end{proof}

\section{DEPT}

\begin{definition}[DEPT]\label{def:b2:structure-determination:dept}
El \emph{DEPT}\index{DEPT} (siglas inglesas de refuerzo sin distorsión por transferencia de polarización) es un conjunto de experimentos
de carbono 13 que dan a las señales de los carbonos con protones unidos un signo o una ausencia según
su número de hidrógenos. Un \emph{carbono cuaternario}\index{carbono cuaternario} no lleva hidrógeno
(está unido a otros cuatro átomos, o es un carbono carbonílico o \omterm{def:b2:huckel:aromatic}{aromático} sin H); no da señal en DEPT.
\end{definition}

\begin{proposition}[Leer un DEPT]\label{prop:b2:structure-determination:dept}
En DEPT-135, los carbonos \ce{CH} y \ce{CH3} apuntan hacia arriba, los carbonos \ce{CH2} hacia abajo y los \omterm{def:b2:structure-determination:dept}{carbonos cuaternarios}
están ausentes. En DEPT-90 solo aparecen los carbonos \ce{CH}.
\end{proposition}

\admitted

El experimento transfiere magnetización de los protones al carbono al que están unidos, mediante una
secuencia de pulsos de radiofrecuencia cuyo último ángulo selecciona la respuesta; cómo lo hacen los pulsos se
explica en el volumen del tercer año. Comparar el espectro desacoplado con el DEPT-135 y el DEPT-90 clasifica cada
línea como C, CH, \ce{CH2} o \ce{CH3}.

\begin{omfigure}
\pgfplotsset{dept/.style={width=\linewidth, height=2.9cm, ycomb, x dir=reverse, ymin=-1, ymax=1.15,
  ytick=\empty, axis y line=none, axis x line=middle, x axis line style={-}, tick label style={font=\tiny},
  title style={font=\scriptsize, at={(0.0,0.95)}, anchor=north west}, every axis plot/.append style={thick}}}
\begin{minipage}[t]{0.49\linewidth}\centering
{\small butan-2-ona}\par
\begin{tikzpicture}
\begin{axis}[dept, xmin=0, xmax=220, xtick={0,50,100,150,200}, title={desacoplado}]
\addplot[omDef] table[x=shift, y=dec] {figdata/out/bachelor-2-nmr-spectra-a.dat};
\end{axis}
\end{tikzpicture}\par
\begin{tikzpicture}
\begin{axis}[dept, xmin=0, xmax=220, xtick={0,50,100,150,200}, title={DEPT-135}]
\addplot[omThm] table[x=shift, y=d135] {figdata/out/bachelor-2-nmr-spectra-a.dat};
\end{axis}
\end{tikzpicture}\par
\begin{tikzpicture}
\begin{axis}[dept, xmin=0, xmax=220, xtick={0,50,100,150,200}, title={DEPT-90}]
\addplot[omMeth] table[x=shift, y=d90] {figdata/out/bachelor-2-nmr-spectra-a.dat};
\end{axis}
\end{tikzpicture}
\end{minipage}\hfill
\begin{minipage}[t]{0.49\linewidth}\centering
{\small benzoato de etilo}\par
\begin{tikzpicture}
\begin{axis}[dept, xmin=0, xmax=220, xtick={0,50,100,150,200}, title={desacoplado}]
\addplot[omDef] table[x=shift, y=dec] {figdata/out/bachelor-2-nmr-spectra-b.dat};
\end{axis}
\end{tikzpicture}\par
\begin{tikzpicture}
\begin{axis}[dept, xmin=0, xmax=220, xtick={0,50,100,150,200}, title={DEPT-135}]
\addplot[omThm] table[x=shift, y=d135] {figdata/out/bachelor-2-nmr-spectra-b.dat};
\end{axis}
\end{tikzpicture}\par
\begin{tikzpicture}
\begin{axis}[dept, xmin=0, xmax=220, xtick={0,50,100,150,200}, title={DEPT-90}]
\addplot[omMeth] table[x=shift, y=d90] {figdata/out/bachelor-2-nmr-spectra-b.dat};
\end{axis}
\end{tikzpicture}
\end{minipage}

{\small Espectros de \ce{^{13}C} desacoplados (desplazamientos y alturas medidos, datos de PubChem) con las respuestas
\omterm{def:b2:structure-determination:dept}{DEPT} (signos según la \cref{prop:b2:structure-determination:dept}, alturas esquemáticas). Butan-2-ona: C=O
209.3 (ausente en \omterm{def:b2:structure-determination:dept}{DEPT}), \ce{CH2} 36.9 (hacia abajo), \ce{CH3} 29.4 y 7.9 (hacia arriba). Benzoato de etilo: el C=O 166.5 y
el carbono del anillo que lleva el éster, 130.6, desaparecen; tres CH \omterm{def:b2:huckel:aromatic}{aromáticos} (132.8, 129.6, 128.3) permanecen en
DEPT-90; el \ce{OCH2} 60.9 apunta hacia abajo.}
% figdata: bachelor-2/nmr-spectra.py parts a, b
\end{omfigure}

\section{Combinar las técnicas}

Cada técnica responde a su propia pregunta, y las respuestas se ensamblan en un orden fijo: primero la fórmula,
porque acota todo lo demás; después los grupos funcionales y el esqueleto; después las
conexiones.

\begin{method}[Resolver una sustancia desconocida]\label{met:b2:structure-determination:solve}
\begin{enumerate}
  \item Fórmula: \omterm{def:b2:mass-spec-atomic:molecular-ion}{ion molecular}, $M+1$ para el número de carbonos, \omterm{def:b2:mass-spec-atomic:isotope-pattern}{perfiles isotópicos}, \omterm{def:b2:mass-spec-atomic:nitrogen-rule}{regla del nitrógeno}; una
        masa exacta si se dispone de ella. Grado de insaturación.
  \item Infrarrojo: \ce{C=O} (y de qué tipo), \ce{O-H}, \ce{N-H}, \ce{C#N}, \ce{C-H} \omterm{def:b2:huckel:aromatic}{aromático} o
        alquénico por encima de \qty{3000}{cm^{-1}}.
  \item Carbono 13 y \omterm{def:b2:structure-determination:dept}{DEPT}: número de tipos de carbono (simetría), sus clases según los desplazamientos, y
        C, CH, \ce{CH2}, \ce{CH3} según el \omterm{def:b2:structure-determination:dept}{DEPT}. Compárese con la fórmula.
  \item RMN de protón: integraciones, desplazamientos, multiplicidades; vecinos según la regla $n + 1$.
  \item Ensámblense los fragmentos; escríbanse todas las estructuras candidatas.
  \item Compruébese cada candidata frente a cada dato, incluidos los fragmentos del \omterm{def:b2:mass-spec-atomic:spectrum}{espectro de masas}; consérvese la que
        los explica todos.
\end{enumerate}
\end{method}

\begin{omfigure}
\begin{tikzpicture}[x=1cm, y=1cm, st/.style={draw=black!70, rounded corners=2pt, fill=white,
  font=\small, align=center, minimum height=0.9cm, inner sep=3pt, text width=2.3cm}]
\node[st, fill=omDef!10] (ms) at (0,0) {EM\\fórmula, $M+1$};
\node[st] (dbe) at (2.9,0) {grado de\\insaturación};
\node[st, fill=omDef!10] (ir) at (5.8,0) {IR\\grupos funcionales};
\node[st, fill=omDef!10] (c13) at (8.7,0) {\ce{^{13}C}, DEPT\\esqueleto};
\node[st, fill=omDef!10] (h1) at (8.7,-1.8) {\ce{^1H}\\conexiones};
\node[st] (cand) at (5.8,-1.8) {estructuras\\candidatas};
\node[st, fill=omThm!10] (chk) at (2.9,-1.8) {comprobar\\cada dato};
\node[st] (ans) at (0,-1.8) {estructura};
\draw[->, thick] (ms) -- (dbe); \draw[->, thick] (dbe) -- (ir); \draw[->, thick] (ir) -- (c13);
\draw[->, thick] (c13) -- (h1); \draw[->, thick] (h1) -- (cand); \draw[->, thick] (cand) -- (chk);
\draw[->, thick] (chk) -- (ans);
\draw[->, thick, densely dashed] (chk.south) -- ++(0,-0.5) -| node[pos=0.25, below, font=\footnotesize]
  {un dato sin explicar: volver a las candidatas} (cand.south);
\end{tikzpicture}
\par\vspace{4pt}

{\small El orden de trabajo ante una sustancia desconocida (\cref{met:b2:structure-determination:solve}). La última etapa
es la que más a menudo se omite: una estructura que deja un pico sin explicar no es todavía la respuesta.}
\end{omfigure}

\section{Casos resueltos}

\begin{example}[Una cetona, C$_5$H$_{10}$O]\label{ex:b2:structure-determination:ketone}
Una sustancia desconocida tiene $M = 86$, una banda IR hacia \qty{1715}{cm^{-1}}, líneas de carbono en 208.9, 45.7, 29.8, 17.4 y
13.7~ppm, y señales de protón en 2.41 (t, 2H), 2.13 (s, 3H), 1.61 (sextuplete, 2H) y 0.91 (t, 3H).
Un grado de insaturación, una cetona saturada (sin protón aldehídico hacia 9.7). Cinco carbonos, cinco líneas:
ninguna simetría. El singlete de tres protones en 2.13 es un \ce{CH3} contiguo al carbonilo; la secuencia triplete,
sextuplete, triplete es un grupo propilo cuyo \ce{CH2} en 2.41 toca el carbonilo. Pentan-2-ona,
\ce{CH3COCH2CH2CH3}. Su isómero simétrico, la pentan-3-ona, mostraría tres líneas de carbono y ningún singlete;
su \omterm{def:b2:mass-spec-atomic:spectrum}{espectro de masas} tiene el ion de McLafferty en 58 que le falta a la pentan-3-ona (\cref{ch:b2:mass-spec-atomic}).
% ledger: c13:pentan-2-one, nmr:pentan-2-one, ir:CO-ketone, ms:pentan-2-one
\end{example}

\begin{example}[Un éster aromático, C$_9$H$_{10}$O$_2$]\label{ex:b2:structure-determination:ester}
Una sustancia desconocida de fórmula \ce{C9H10O2} (cinco grados de insaturación) muestra una banda de carbonilo y siete
líneas de carbono: 166.5 (C), 132.8 (CH), 130.6 (C), 129.6 (CH), 128.3 (CH), 60.9 (\ce{CH2}) y 14.3
(\ce{CH3}). Su espectro de protón: 8.05 (m, 2H), 7.35--7.63 (m, 3H), 4.37 (c, 2H), 1.38 (t, 3H). Un
anillo bencénico monosustituido (cuatro líneas de carbono \omterm{def:b2:huckel:aromatic}{aromático}, dos de ellas de dos carbonos cada una; cinco
protones \omterm{def:b2:huckel:aromatic}{aromáticos}) da cuenta de cuatro insaturaciones, y el carbonilo, de la quinta. El par cuadruplete--triplete
es un grupo etilo; su \ce{CH2} en 4.37 (y en 60.9 en carbono) está unido a un oxígeno. El carbonilo de éster
en 166.5 está unido al anillo: los dos protones en 8.05, contiguos a él, están desplazados a campo bajo. Benzoato de
etilo, \ce{C6H5COOCH2CH3}.
% ledger: c13:ethylbenzoate, nmr:ethylbenzoate
\end{example}

\section{Trampas}

\begin{itemize}
  \item \emph{Protones intercambiables} (\ce{O-H}, \ce{N-H}): dan señales anchas a desplazamientos variables, a
        menudo no están acoplados a sus vecinos, y desaparecen cuando se agita la muestra con una gota de \ce{D2O}.
  \item \emph{Señales solapadas}: dos carbonos de un anillo bencénico, o varios \ce{CH2} de una cadena, pueden
        compartir una línea; cuéntense las líneas frente a la fórmula antes de concluir sobre la simetría.
  \item \emph{Protones diastereotópicos}: los dos protones de un \ce{CH2} contiguo a un estereocentro no son
        equivalentes, pueden tener desplazamientos distintos y acoplarse entre sí.
  \item \emph{Sistemas de segundo orden}: cuando dos protones acoplados tienen desplazamientos más próximos que unas pocas veces su
        constante de acoplamiento, los multipletes se inclinan uno hacia otro y la regla $n+1$ falla; un
        espectrómetro de campo más intenso restablece multipletes sencillos.
\end{itemize}

\begin{safety}{\ghs[5mm]{GHS02}\,\ghs[5mm]{GHS07}\,\ghs[5mm]{GHS08}}
La butan-2-ona es inflamable. El imán de un espectrómetro de RMN está siempre activo: nada de herramientas de acero, botellas de
gas ni tarjetas magnéticas cerca de él, y ningún acceso para personas con marcapasos.
% ledger: ghs:butanone, ghs:benzylacetate
\end{safety}

\section{Ejercicios}

\begin{exercise}[$\star$]\label{exo:b2:structure-determination:1}
¿Cuántas líneas muestra el espectro de \ce{^{13}C} desacoplado de cada xileno (1,2-, 1,3- y
1,4-dimetilbenceno)?
\end{exercise}

\begin{exercise}[$\star$]\label{exo:b2:structure-determination:2}
Prevéanse los espectros DEPT-135 y DEPT-90 del butan-2-ol.
\end{exercise}

\begin{exercise}[$\star$]\label{exo:b2:structure-determination:3}
¿A qué tipo de carbono pertenece con más probabilidad una línea en 209, 170, 129, 64 y 14~ppm?
\end{exercise}

\begin{exercise}[$\star$]\label{exo:b2:structure-determination:4}
Calcúlese el grado de insaturación de \ce{C8H8O2} y propóngase una estructura que contenga un anillo bencénico.
\end{exercise}

\begin{exercise}[$\star\star$]\label{exo:b2:structure-determination:5}
Un líquido \ce{C4H8O} muestra una banda IR en \qty{1715}{cm^{-1}} y líneas de carbono en 209, 37, 29 y 8~ppm
(DEPT-135: 37 hacia abajo, 29 y 8 hacia arriba, 209 ausente). Identifíquese.
\end{exercise}

\begin{exercise}[$\star\star$]\label{exo:b2:structure-determination:6}
¿Cómo se distinguiría la pentan-2-ona de la pentan-3-ona solo con RMN de \ce{^{13}C}? ¿Con RMN de protón? ¿Con
\omterm{def:b2:mass-spec-atomic:spectrum}{espectrometría de masas}?
\end{exercise}

\begin{exercise}[$\star\star$]\label{exo:b2:structure-determination:7}
Un éster tiene $M = 88$ y señales de protón en 4.12 (c, 2H), 2.04 (s, 3H) y 1.26 (t, 3H) (datos de
ejercicio). Identifíquese.
\end{exercise}

\begin{exercise}[$\star\star$]\label{exo:b2:structure-determination:8}
Distíngase el 1,2-diclorobenceno del 1,4-diclorobenceno por RMN de \ce{^{13}C}.
\end{exercise}

\begin{exercise}[$\star\star$]\label{exo:b2:structure-determination:9}
En el 2-clorobutano, ¿por qué no son equivalentes los dos protones del grupo \ce{CH2}? ¿Qué efecto tiene eso sobre
el espectro de protón?
\end{exercise}

\begin{exercise}[$\star\star\star$]\label{exo:b2:structure-determination:10}
Una sustancia desconocida \ce{C9H10O} muestra una banda IR en \qty{1690}{cm^{-1}}, líneas de carbono en 200.8 (C), 137.0 (C),
132.9 (CH), 128.6 (CH), 128.0 (CH), 31.8 (\ce{CH2}) y 8.3 (\ce{CH3}), y señales de protón en 7.95 (m,
2H), 7.40--7.55 (m, 3H), 2.98 (c, 2H) y 1.22 (t, 3H) (datos de ejercicio). Identifíquese, y explíquese el
bajo número de onda del carbonilo.
\end{exercise}

\begin{exercise}[$\star\star\star$]\label{exo:b2:structure-determination:11}
Cuatro isómeros \ce{C5H10O2}: ácido pentanoico, butanoato de metilo, propanoato de etilo y etanoato de propilo.
Dese para cada uno un rasgo de su espectro IR o de RMN de protón que lo distinga de los otros tres.
\end{exercise}

\begin{exercise}[$\star\star\star$]\label{exo:b2:structure-determination:12}
Un informe asigna la línea en 166.5~ppm del benzoato de etilo al carbono del anillo que lleva el éster y
la línea en 130.6~ppm al carbono carbonílico. Muéstrese, con los datos \omterm{def:b2:structure-determination:dept}{DEPT} y el gráfico de desplazamientos, que la
asignación no puede ser correcta.
\end{exercise}

\section{Problema: cuatro espectros del jazmín}

\begin{problem}[{Problema de fin de semana --- la fórmula a partir del ion molecular, la banda infrarroja del carbonilo,
el carbono 13 y el DEPT, el espectro de protón, y una estructura contrastada con los fragmentos}]
\label{pb:b2:structure-determination:1}
Un líquido incoloro de olor floral a jazmín da los datos siguientes. \omterm{def:b2:mass-spec-atomic:spectrum}{Espectro de masas} (NIST
WebBook): $m/z$ 150 (31.7), 151 (3.1), 108 (100), 91 (71.7), 90 (40.6), 43 (37.6). Infrarrojo (datos de
ejercicio): bandas intensas en 1740 y \qty{1230}{cm^{-1}}, bandas débiles justo por encima de \qty{3000}{cm^{-1}}, ninguna
banda ancha hacia \qty{3300}{cm^{-1}}. Carbono 13, CDCl$_3$ (datos de PubChem): 170.70, 136.14, 128.56,
128.24, 66.24, 20.82~ppm. Protón, CDCl$_3$, \qty{90}{MHz}: 7.33 (m, 5H), 5.09 (s, 2H), 2.06 (s, 3H).
% ledger: use:benzylacetate, ms:benzylacetate, c13:benzylacetate, nmr:benzylacetate, ir:CO-ester, iso:C12, iso:C13, mass:C12, mass:H1, mass:O16

\begin{center}
\begin{tikzpicture}
\pgfplotsset{dept/.style={width=0.8\linewidth, height=3.2cm, ycomb, x dir=reverse, ymin=0, ymax=1.15,
  xmin=0, xmax=200, ytick=\empty, axis y line=none, axis x line=bottom, x axis line style={-}, tick label style={font=\tiny},
  title style={font=\scriptsize, at={(0.02,0.82)}, anchor=west}, every axis plot/.append style={thick}}}
\begin{axis}[dept, title={desacoplado}, xtick={0,20,40,60,80,100,120,140,160,180,200}]
\addplot[omDef] table[x=shift, y=dec] {figdata/out/bachelor-2-nmr-spectra-c.dat};
\end{axis}
\end{tikzpicture}
\end{center}
% figdata: bachelor-2/nmr-spectra.py part c

\medskip
\noindent\textbf{Parte I --- La fórmula.}
\begin{enumerate}
  \item A partir de la razón 151/150, estímese el número de carbonos.
  \item ¿Qué dice la \omterm{def:b2:mass-spec-atomic:nitrogen-rule}{regla del nitrógeno}?
  \item Con nueve carbonos, ¿qué queda de la masa 150? Propóngase una fórmula con oxígeno.
  \item ¿Por qué es menos probable \ce{C10H14O}, también de masa nominal 150?
  \item Calcúlese el grado de insaturación.
  \item Calcúlese la \omterm{def:b2:mass-spec-atomic:exact-mass}{masa monoisotópica} de la fórmula elegida.
\end{enumerate}

\medskip
\noindent\textbf{Parte II --- Infrarrojo.}
\begin{enumerate}[resume]
  \item Asígnese la banda en \qty{1740}{cm^{-1}}. ¿Qué familia sugiere su posición?
  \item Asígnese la banda en \qty{1230}{cm^{-1}}.
  \item ¿Qué indican las bandas débiles por encima de \qty{3000}{cm^{-1}}?
  \item ¿Qué excluye la ausencia de una banda ancha hacia \qty{3300}{cm^{-1}}?
  \item ¿Está el carbonilo conjugado con el anillo? Explíquese a partir de su número de onda.
\end{enumerate}

\medskip
\noindent\textbf{Parte III --- Carbono 13 y \omterm{def:b2:structure-determination:dept}{DEPT}.}
\begin{enumerate}[resume]
  \item Asígnese cada línea a un tipo de carbono.
  \item ¿Qué línea apunta hacia abajo en DEPT-135?
  \item ¿Qué líneas desaparecen en DEPT-135?
  \item ¿Qué líneas permanecen en DEPT-90?
  \item ¿Cuántos tipos de carbono tiene un anillo bencénico monosustituido? Explíquese por qué solo aparecen dos líneas
        de CH \omterm{def:b2:huckel:aromatic}{aromático}.
  \item La línea en 128.24 es la más alta. ¿Representa el mayor número de carbonos?
\end{enumerate}

\medskip
\noindent\textbf{Parte IV --- RMN de protón y estructura.}
\begin{enumerate}[resume]
  \item Asígnense las tres señales de protón.
  \item ¿Por qué son singletes las señales en 5.09 y 2.06?
  \item ¿Por qué está tan desplazado a campo bajo el \ce{CH2} en 5.09?
  \item Ensámblese la estructura y nómbrese.
  \item Su isómero, el 2-feniletanoato de metilo, \ce{C6H5CH2COOCH3}, tiene la misma fórmula. ¿Qué dato lo
        descarta?
  \item Explíquense los fragmentos en 108, 91 y 43.
  \item ¿Cuántas líneas mostraría un espectro de \ce{^{13}C} desacoplado perfectamente resuelto del compuesto, y
        cuántas de ellas apuntan hacia abajo en DEPT-135?
\end{enumerate}
\end{problem}
